%% Question bank chapter file - generated by bank_tools.py build (do not edit while the bank is still being built:
%% it is regenerated from the per-test data in Question Bank\data). Every question is one bankq block:
%%   \begin{bankq}[id=..., source={...}, paper=.., qno=.., type=..., subtopic={...}, answer={...}, answer_src=key|solved,
%%                 working={...}, difficulty=1-10, diff_parts={K S I C F}, diff_note={...},
%%                 group=<passage id or empty>, also_in={duplicate copies in other tests}]
%% difficulty = min(10, K+S+I+C+F-1): K concepts 1-3, S steps 1-4, I insight 0-2, C computation 0-2, F format 0-1
%%   ...question LaTeX, options, figures (\bfig / \bfigi -> img/)...
%%   \end{bankq}
%% Passages shared by several questions are bankpassage blocks; their questions carry group=<passage id>.
\documentclass[10pt]{article}
\usepackage[a4paper,margin=1.8cm]{geometry}
\usepackage[T1]{fontenc}
\usepackage{amsmath,amssymb,mathtools}
\usepackage{graphicx}
\usepackage[export]{adjustbox}
\usepackage{xcolor,needspace,keyval}
\usepackage{enumitem}
\usepackage{booktabs,array,tabularx}
\usepackage[hidelinks]{hyperref}
\setlength{\parindent}{0pt}
\setlength{\parskip}{3pt}
\makeatletter
\define@key{bq}{id}{\def\bq@id{#1}}\define@key{bq}{source}{\def\bq@src{#1}}\define@key{bq}{paper}{\def\bq@paper{#1}}
\define@key{bq}{qno}{\def\bq@qno{#1}}\define@key{bq}{type}{\def\bq@type{#1}}\define@key{bq}{subtopic}{\def\bq@sub{#1}}
\define@key{bq}{answer}{\def\bq@ans{#1}}\define@key{bq}{answer_src}{\def\bq@asrc{#1}}\define@key{bq}{working}{\def\bq@work{#1}}
\define@key{bq}{group}{\def\bq@group{#1}}\define@key{bq}{also_in}{\def\bq@also{#1}}
\define@key{bq}{difficulty}{\def\bq@diff{#1}}\define@key{bq}{diff_parts}{\def\bq@dparts{#1}}\define@key{bq}{diff_note}{\def\bq@dnote{#1}}
\newenvironment{bankq}[1][]{\def\bq@work{}\def\bq@also{}\def\bq@group{}\def\bq@diff{?}\def\bq@dparts{}\def\bq@dnote{}\setkeys{bq}{#1}%
  \par\needspace{4\baselineskip}\medskip\noindent{\small\textbf{\bq@id}\hfill\color{gray}[\bq@src, Paper \bq@paper, Q\bq@qno]
  \bq@type\ $\cdot$ \bq@sub\ $\cdot$ \textbf{D\bq@diff}}\par\nopagebreak\relax}
  {\par{\small\color{teal}Answer: \textbf{\bq@ans}\ifx\bq@work\@empty\else\ (solved: \bq@work)\fi
  \ifx\bq@also\@empty\else\par\color{gray}Also in: \bq@also\fi}\par}
\define@key{bp}{id}{\def\bp@id{#1}}\define@key{bp}{source}{\def\bp@src{#1}}\define@key{bp}{questions}{\def\bp@qs{#1}}
\newenvironment{bankpassage}[1][]{\setkeys{bp}{#1}\par\needspace{4\baselineskip}\bigskip\noindent
  {\small\textbf{Passage \bp@id}\hfill\color{gray}[\bp@src; questions \bp@qs]}\par\nopagebreak\relax\itshape}{\par}
\makeatother
\newcommand{\bfig}[2]{\par\smallskip{\centering\includegraphics[width=#1pt,max width=0.62\linewidth,max height=0.32\textheight]{img/#2}\par}\smallskip}
\newcommand{\bfigi}[2]{\includegraphics[width=#1pt,max width=\linewidth,max height=0.32\textheight,valign=c]{img/#2}}
\newcommand{\bmissing}[1]{\textit{[figure missing in the source conversion: \detokenize{#1}]}}
\newenvironment{qoptions}{\begin{itemize}[leftmargin=1.8em,itemsep=1pt,topsep=2pt,label={}]}{\end{itemize}}
\newcommand{\opt}[2]{\item[(#1)] #2}
\providecommand{\fwstart}{}\providecommand{\fwend}{}\providecommand{\qfullsplit}{}
\newcolumntype{L}[1]{>{\raggedright\arraybackslash}p{#1}}
\newcolumntype{C}[1]{>{\centering\arraybackslash}p{#1}}
\newcolumntype{R}[1]{>{\raggedleft\arraybackslash}p{#1}}
\providecommand{\sectionhdr}[1]{}\providecommand{\subsectionhdr}[1]{}


\begin{document}
{\Large\bfseries AOD --- Question Bank}\par\smallskip
{\small Mathematics chapter M24. Questions from all converted test papers, grouped by source series. Answers: printed key where it matches the paper, otherwise solved (answer\_src=solved, with working). Exact repeats across papers appear once, with also\_in listing the other copies.}\par\medskip\hrule

\section*{Score-I papers}
\begin{bankq}[id=SI12-P2-MATH-S2-Q05, source={Score-I - Test 12 (09 Mar 2022)}, paper=2, qno=5, type=single, subtopic={Mean value theorems}, answer={0.00}, answer_src=key, difficulty=4, diff_parts={K1 S2 I2 C0 F0}, diff_note={concepts: mean value theorem bound | steps: bound the maximum of the modulus of f on a short first interval by 2 times its length times that maximum; conclude f = 0 there; repeat step by step up to x = 5 | I: iterative MVT bound}]
Let $f$ be a differentiable function on $[1,5]$, $f(1)=0$, $|f'(x)|\le2|f(x)|$ on $[1,5]$, then the value of $\displaystyle\sum_{r=1}^5f(r)$ is
\end{bankq}

\section*{TOAS papers}
\begin{bankq}[id=TOAS06-P1-MATH-S2-Q11, source={TOAS - Test 06 (17 Apr 2022)}, paper=1, qno=11, type=multi, subtopic={Mean value theorems}, answer={A,C}, answer_src=key, difficulty=5, diff_parts={K2 S2 I1 C0 F1}, diff_note={concepts: Rolle's theorem, intermediate value | steps: roots between consecutive poles; count; chord theorem on [0,2] | I: continuity of f(x + 1) - f(x)}]
Which of the following statements is/are always true ?
\begin{qoptions}
\opt{A}{number of roots of equation $\displaystyle\sum_{r=1}^{2017}\frac{r}{x-r}=0$ is 2016}
\opt{B}{number of roots of equation $\displaystyle\sum_{r=1}^{2017}\frac{r}{x-r}=0$ is 2017}
\opt{C}{if $f(x)$ is continuous in $[0,2]$ such that $f(0)=f(2)$, then there exists atleast one $x\in[0,1]$ such that $f(x)=f(x+1)$.}
\opt{D}{if $f(x)$ is continuous in $[0,2]$ such that $f(0)=f(2)$, then there exists atleast one $x\in(0,1)$ such that $f(x)=f(x+1)$.}
\end{qoptions}
\end{bankq}
\begin{bankq}[id=TOAS07-P1-MATH-S4-Q08, source={TOAS - Test 07 (01 May 2022)}, paper=1, qno=8, type=numerical, subtopic={Rate of change}, answer={12}, answer_src=key, difficulty=5, diff_parts={K2 S2 I0 C1 F1}, diff_note={concepts: tangent intercepts, related rates | steps: position at t = 2; intercepts of the tangent in terms of x; area A = x^(3/2)/4; differentiate and substitute}]
If a particle starts moving from $(0,0)$ and rate of change of its abscissa is 8 cm/s on the curve $y=\sqrt x$, then rate of change of area of triangle formed by co-ordinate axes \& tangent at a point reached after 2 sec of start is
\end{bankq}
\begin{bankpassage}[id=TOAS07-P2-MATH-S3-G3-4, source={TOAS - Test 07 (01 May 2022)}, questions={3-4}]
$A(0,0)$ and $B(8,2)$ are fixed points on the curve $y^3=|x|$. A point $C$ starts moving from origin along the curve $(x<0)$ such that rate of change in the ordinate is 2cm/sec. After $t_0$ seconds triangle $ABC$ become a right triangle.
\end{bankpassage}
\begin{bankq}[id=TOAS07-P2-MATH-S3-Q03, source={TOAS - Test 07 (01 May 2022)}, paper=2, qno=3, type=numerical, subtopic={Rate of change}, answer={19.00}, answer_src=key, difficulty=5, diff_parts={K2 S2 I0 C1 F1}, diff_note={concepts: related rates, tangent and area | steps: position after t0; area of the triangle from the tangent; differentiate and substitute}, group=TOAS07-P2-MATH-S3-G3-4]
Rate of change in area of trinagle after $t_0$ second is $k$ then $2k$ is equal to
\end{bankq}
\begin{bankq}[id=TOAS07-P2-MATH-S3-Q04, source={TOAS - Test 07 (01 May 2022)}, paper=2, qno=4, type=numerical, subtopic={Rate of change}, answer={1.00}, answer_src=key, difficulty=3, diff_parts={K1 S2 I0 C0 F1}, diff_note={concepts: tangent meeting a curve again | steps: slope at C; tangent line; second intersection point}, group=TOAS07-P2-MATH-S3-G3-4]
After $t_0$ seconds, tangent is drawn to the curve at point $C$ to intersect it again at $(\alpha,\beta)$ then $4\alpha+3\beta$ is -
\end{bankq}
\begin{bankq}[id=TOAS08-P1-MATH-S1-Q01, source={TOAS - Test 08 (15 May 2022)}, paper=1, qno=1, type=multi, subtopic={Mean value theorems}, answer={A,B,C}, answer_src=key, difficulty=5, diff_parts={K1 S2 I2 C0 F1}, diff_note={concepts: Rolle's theorem with an antiderivative | steps: find Q(x) with Q' = P(x); evaluate Q at 0, 5^(1/5), 3^(1/3); apply Rolle to the intervals | I: construct the antiderivative}]
Consider $P(x)=9x^8-18x^5-20x^3+15$, then $P(x)=0$ has
\begin{qoptions}
\opt{A}{atleast one real root in $[0,3^{1/3}]$}
\opt{B}{atleast one real root in $[5^{1/5},3^{1/3}]$}
\opt{C}{atleast one real root in $[0,5^{1/5}]$}
\opt{D}{no real root in $[0,3^{1/3}]$}
\end{qoptions}
\end{bankq}
\begin{bankpassage}[id=TOAS08-P1-MATH-S2-G7-8, source={TOAS - Test 08 (15 May 2022)}, questions={7-8}]
Let $y=f(x)$ be a fourth degree algebraic polynomial function such that\newline
(i) $\displaystyle\lim_{x\to1}\dfrac{f(x)}{(x-1)^2}=2$\newline
(ii) It cuts $y$-axis a point where slope of normal is $\dfrac18$\newline
(iii) Tangent at $x=2$ is parallel to $y=8x+1$\newline
Consider $y=g(x)$ be a quadratic curve touching and cutting $y=f(x)$ on $x$-axis and $y$-axis respectively.
\end{bankpassage}
\begin{bankq}[id=TOAS08-P1-MATH-S2-Q07, source={TOAS - Test 08 (15 May 2022)}, paper=1, qno=7, type=single, subtopic={Tangents and normals}, answer={B}, answer_src=key, difficulty=1, diff_parts={K1 S1 I0 C0 F0}, diff_note={concepts: angle of intersection of curves | steps: slopes of both curves at the y-axis; tan of the angle between them}, group=TOAS08-P1-MATH-S2-G7-8]
The angle of intersection between $y=f(x)$ \& $y=g(x)$ at point where they cut $y$-axis is -
\begin{qoptions}
\opt{A}{$\tan^{-1}\dfrac{2}{33}$}
\opt{B}{$\tan^{-1}\dfrac{2}{49}$}
\opt{C}{$\tan^{-1}\dfrac{14}{33}$}
\opt{D}{$\tan^{-1}\dfrac{14}{49}$}
\end{qoptions}
\end{bankq}
\begin{bankq}[id=TOAS08-P1-MATH-S2-Q08, source={TOAS - Test 08 (15 May 2022)}, paper=1, qno=8, type=single, subtopic={Tangents and normals}, answer={D}, answer_src=key, difficulty=1, diff_parts={K1 S1 I0 C0 F0}, diff_note={concepts: subtangent | steps: subtangent = y/y'; solve for the abscissa}, group=TOAS08-P1-MATH-S2-G7-8]
The abscissae of point on curve $y=g(x)$, where length of subtangent is 2018 can be -
\begin{qoptions}
\opt{A}{1010}
\opt{B}{2019}
\opt{C}{3028}
\opt{D}{4037}
\end{qoptions}
\end{bankq}
\begin{bankq}[id=TOAS08-P1-MATH-S3-Q01, source={TOAS - Test 08 (15 May 2022)}, paper=1, qno=1, type=numerical, subtopic={Tangents and normals}, answer={2.00}, answer_src=key, difficulty=4, diff_parts={K1 S2 I1 C0 F1}, diff_note={concepts: tangents with given slope on an implicit curve | steps: differentiate implicitly; slope = -1/3; solve for sin(x + 2y); count points in the range | I: slope condition fixes the angle}]
The number of tangents to the curve $y=\cos(x+2y)$, $-2\pi\leqslant x\leqslant2\pi$, that are parallel to the line $x+3y=0$ is
\end{bankq}
\begin{bankq}[id=TOAS08-P1-MATH-S3-Q02, source={TOAS - Test 08 (15 May 2022)}, paper=1, qno=2, type=numerical, subtopic={Mean value theorems}, answer={3.67}, answer_src=solved, working={Roots of $f'(x)(1-f^2)=0$: $f=1$ has 5 forced roots, $f=-1$ has 2 (sign changes between integers), and Rolle gives at least 4 distinct points with $f'=0$: $\lambda=11$, $\lambda/3\approx3.67$}, difficulty=5, diff_parts={K1 S2 I2 C0 F1}, diff_note={concepts: Rolle's theorem and intermediate values | steps: zeros of 1 - f^2 from sign changes; zeros of f' from Rolle's theorem; combine distinct points | I: both factors give roots}]
Let $f(x)$ be a differentiable function on $[0,8]$ such that $f(1)=5$, $f(2)=\dfrac12$, $f(3)=6$, $f(4)=-3$, $f(5)=4$, $f(6)=\dfrac14$ and $f(7)=\dfrac{-1}{5}$. If the minimum number of roots of the equation $f'(x)-f'(x)(f(x))^2=0$ is $\lambda$, then $\dfrac\lambda3$ equals
\end{bankq}
\begin{bankq}[id=TOAS08-P1-MATH-S3-Q07, source={TOAS - Test 08 (15 May 2022)}, paper=1, qno=7, type=numerical, subtopic={Rate of change}, answer={0.53}, answer_src=key, difficulty=4, diff_parts={K1 S2 I0 C1 F1}, diff_note={concepts: related rates in a tank | steps: volume of water in terms of length and height; differentiate; substitute the given values}]
If a rectangular tank has a sliding panel $S$ that divides it into two adjustable tanks of width 3 ft. See figure. Water is poured into the left compartment at the rate of 5 ft$^3$/min. At the same time $S$ is moved to the right at the rate of 3ft/min. When the left comparment is 10 ft long it contains 70ft$^3$ of water, then how fast the water level is falling
\bfig{106.8}{TOAS08_p1_math_s3_q07_1.png}
\end{bankq}
\begin{bankq}[id=TOAS08-P2-MATH-S3-Q07, source={TOAS - Test 08 (15 May 2022)}, paper=2, qno=7, type=numerical, subtopic={Tangents and normals}, answer={8.00}, answer_src=key, difficulty=4, diff_parts={K1 S2 I0 C1 F1}, diff_note={concepts: tangent to a parametric curve | steps: point P at t = pi/4; tangent line; second intersection from the parameter; distance PQ}]
A curve is represented by $x=\sec^2t$ and $y=\cot t$, where $t$ is the parameter. If the tangent at point $P$ on curve where $t=\dfrac\pi4$ meets the curve again at $Q$ then $|PQ|=\dfrac{p\sqrt q}{2}$, where $p,q$ are coprime. Find $p+q$
\end{bankq}
\begin{bankq}[id=TOAS09-P1-MATH-S1-Q03, source={TOAS - Test 09 (30 May 2022)}, paper=1, qno=3, type=single, subtopic={Tangents and normals}, answer={C}, answer_src=key, difficulty=3, diff_parts={K2 S2 I0 C0 F0}, diff_note={concepts: tangent to y^3 = x^4, normal to a circle | steps: normal to the circle passes through its centre (1,0); tangent condition at x0; solve x0 = 4; m and (3m/4)^3}]
If $m$ is the slope of a line which is tangent to $y^3=x^4$ \& a normal to $x^2-2x+y^2=0$, then $\left(\dfrac{3m}{4}\right)^3$ is equal to ($m\neq0$)
\begin{qoptions}
\opt{A}{3}
\opt{B}{$\dfrac43$}
\opt{C}{4}
\opt{D}{$\dfrac34$}
\end{qoptions}
\end{bankq}
\begin{bankq}[id=TOAS09-P1-MATH-S1-Q04, source={TOAS - Test 09 (30 May 2022)}, paper=1, qno=4, type=single, subtopic={Mean value theorems}, answer={C}, answer_src=key, difficulty=4, diff_parts={K1 S2 I2 C0 F0}, diff_note={concepts: Lagrange mean value theorem | steps: define a function whose derivative gives the equation; apply MVT on [a,b]; conclude a root exists | I: construct the function}]
For the equation $b\sqrt{1-a^2}-a\sqrt{1-b^2}=\sin\!\left(\dfrac{b-a}{\sqrt{1-x^2}}\right)$, where $0<a<b<1/2$, then which of the following is true -
\begin{qoptions}
\opt{A}{there is no real root of the equation in $(a,b)$}
\opt{B}{there is exactly one real root of the equation in $(a,b)$}
\opt{C}{there is atleast one root of the equation in $(a,b)$}
\opt{D}{none of these}
\end{qoptions}
\end{bankq}
\begin{bankq}[id=TOAS09-P2-MATH-S1-Q07, source={TOAS - Test 09 (30 May 2022)}, paper=2, qno=7, type=single, subtopic={Tangents and normals}, answer={C}, answer_src=key, difficulty=1, diff_parts={K1 S1 I0 C0 F0}, diff_note={concepts: slope of normal and tangent | steps: slope of the normal tan(2 pi/3) = -sqrt3 so f'(4) = 1/sqrt3}]
If the normal to the curve $y=f(x)$ at the point $(4,6)$ makes an angle $\dfrac{2\pi}{3}$ with the positive $x$ axis in anticlockwise direction, then $f'(4)$ is
\begin{qoptions}
\opt{A}{$\sqrt3$}
\opt{B}{$-\sqrt3$}
\opt{C}{$\dfrac{1}{\sqrt3}$}
\opt{D}{$\dfrac{-1}{\sqrt3}$}
\end{qoptions}
\end{bankq}
\begin{bankq}[id=TOAS09-P2-MATH-S2-Q08, source={TOAS - Test 09 (30 May 2022)}, paper=2, qno=8, type=numerical, subtopic={Tangents and normals}, answer={17}, answer_src=solved, working={$y=x^4+1$ from the data; at $Q$: $y'=\frac12$, normal length $=\frac{17}{16}\sqrt{1+\frac14}=\frac{17\sqrt5}{32}$}, difficulty=4, diff_parts={K2 S2 I0 C0 F1}, diff_note={concepts: differential equation, normal length | steps: y' = c x^3 integrates to y = c x^4/4 + 1; use Q to fix c; normal length at Q}]
Let $y=f(x)$ be a curve passing through $P(0,1)$ \& $Q\left(\dfrac12,\dfrac{17}{16}\right)$ and having a property that slope of the tangent at any point is proportional to the cube of the abscissa at the point. If length of normal at Q be $\ell$ then $\dfrac{32\ell}{\sqrt5}$ is equal to
\end{bankq}
\end{document}
